\subsection{Odd-order ordinary shadows at a single \texorpdfstring{$A_6$}{A6} chief step}

\begin{statement}{D.33}{Proved}{Every odd-order multiplier-one ordinary shadow lifts across a single $A_6$ kernel}
Let
\[
 1\longrightarrow S\longrightarrow E
       \xrightarrow{\pi}D\longrightarrow1,
 \qquad S\simeq A_6,
 \tag{D.33.1}
\]
be a finite extension, with a given epimorphism
$q_E:\widehat F\twoheadrightarrow E$ and
peripheral values $X_i$, whose images in $D$
are $d_i$.  Let $\alpha\in\operatorname{Aut}(D)$
have odd order $m$ and satisfy
$\alpha(d_i)\in d_i^D$ for the three labelled
peripherals.  There is an actual automorphism
$\widetilde\alpha$ of $E$, of order $m$,
covering $\alpha$, fixing $S$ pointwise, and
satisfying
\[
                    \widetilde\alpha(X_i)\in X_i^E.
 \tag{D.33.2}
\]
Consequently $\widetilde\alpha q_E$ is an
ordinary peripheral epimorphic lift of
$\alpha\pi q_E$.  In particular a
multiplier-one ordinary obstruction of
order five cannot occur at a single
$A_6$ kernel.
\end{statement}

\begin{proof}
Use the actual identification
$\operatorname{Aut}(S)=\operatorname{Aut}(A_6)$
and $\operatorname{Out}(S)=O=\mathbf F_2^2$
proved by the complete automorphism
calculation in D.31.
The extension determines a homomorphism
$\epsilon:D\to O$.
Centerlessness gives the actual isomorphism
\[
 E\xrightarrow{\sim}
       \operatorname{Aut}(S)\times_{O,\epsilon}D,
 \qquad e\longmapsto(\operatorname{Ad}_e,\pi(e)).
 \tag{D.33.3}
\]
Its kernel is $Z(S)=1$, and any compatible
automorphism over a given $d$ differs from
the conjugation action of a chosen lift
of $d$ by an inner automorphism of $S$.
This proves surjectivity as well.

Choose $h_i\in D$ with
$\alpha(d_i)=h_i d_i h_i^{-1}$.
Since $O$ is abelian,
$\epsilon(\alpha(d_i))=\epsilon(d_i)$.
The elements $d_1,d_2$ generate $D$,
so the two homomorphisms agree:
\[
                         \epsilon\alpha=\epsilon.
 \tag{D.33.4}
\]
Define the actual cumulative conjugator
\[
 c_i=\alpha^{m-1}(h_i)\,
          \alpha^{m-2}(h_i)\cdots\alpha(h_i)h_i.
 \tag{D.33.5}
\]
Iterating the conjugacy equation gives
$\alpha^m(d_i)=c_i d_i c_i^{-1}$.
Thus $c_i\in C_D(d_i)$.
Equation \textup{(D.33.4)} and the
oddness of $m$ give
\[
                 \epsilon(c_i)
                  =m\epsilon(h_i)=\epsilon(h_i).
 \tag{D.33.6}
\]
Put $h_i'=h_i c_i^{-1}$.
The centralizer membership proves
$h_i'd_i(h_i')^{-1}=\alpha(d_i)$,
and \textup{(D.33.6)} proves
$\epsilon(h_i')=0$.
This is the exact adjustment allowed
by ordinary peripheral conjugacy.

In the pullback \textup{(D.33.3)},
define
\[
 \widetilde\alpha:(g,d)\longmapsto(g,\alpha(d)).
 \tag{D.33.7}
\]
Equation \textup{(D.33.4)} makes it
a homomorphism of that very group,
with inverse $(g,d)\mapsto(g,\alpha^{-1}(d))$.
It fixes its kernel pointwise.
It has order exactly $m$: its
$m$th power is the identity, and
its projection to the base has
order $m$.
If $X_i=(g_i,d_i)$, then
$(1,h_i')$ belongs to the pullback
because $\epsilon(h_i')=0$, and
\[
 (1,h_i')\,X_i\,(1,h_i')^{-1}
     =(g_i,\alpha(d_i))
     =\widetilde\alpha(X_i).
 \tag{D.33.8}
\]
Transporting \textup{(D.33.7)--(D.33.8)}
back through the actual isomorphism
\textup{(D.33.3)} proves the statement.
\end{proof}

D.33 concerns the ordinary fibre:
the $h_i'$ need not equal the
initially prescribed $h_i$.
Their difference is exactly the
centralizer element
\textup{(D.33.5)}, so the result
does not erase the labelled
obstruction in D.31.
It also gives a precise limitation
on changing the order-two example
of D.32 into an order-five example
while retaining a single $A_6$
kernel.  For an $A_6$ chief kernel,
an ordinary order-five obstruction
would have to involve more than
one simple factor.
In that case the actual outer
group is
\[
       \operatorname{Out}(A_6^\Omega)
         =(\mathbf F_2^2)^\Omega
                         \rtimes\operatorname{Sym}(\Omega),
 \tag{D.33.9}
\]
with its transitive factor action.
The argument proving
\textup{(D.33.4)} used the abelian
single-factor outer quotient;
the displayed permutation
component must be retained in
an analysis of that remaining case.
